\documentclass[a4paper,11pt,reqno]{amsart}

\usepackage[margin=1in]{geometry}
\usepackage[T1]{fontenc}
\usepackage{amsmath,amssymb,mathtools}
\usepackage{booktabs}
\usepackage{newtxtext,newtxmath}
\usepackage{microtype}
\usepackage[numbers,sort&compress]{natbib}
\usepackage{enumitem}
\usepackage{needspace}
\usepackage{etoolbox}
\usepackage{xurl}
\usepackage[hidelinks,pdfencoding=auto,psdextra]{hyperref}
\usepackage{bookmark}

\setcounter{secnumdepth}{2}
\setcounter{tocdepth}{1}
\setlength{\parindent}{1.5em}
\setlength{\parskip}{0pt}
\setlength{\emergencystretch}{3em}
\setlist[itemize]{leftmargin=2em,itemsep=0.15em,topsep=0.35em}
\setlist[enumerate]{leftmargin=2.2em,itemsep=0.15em,topsep=0.35em}
\allowdisplaybreaks[2]
\raggedbottom
\frenchspacing

\theoremstyle{plain}
\newtheorem{theorem}{Theorem}[section]
\newtheorem{proposition}[theorem]{Proposition}
\newtheorem{lemma}[theorem]{Lemma}
\newtheorem{corollary}[theorem]{Corollary}
\newtheorem{conjecture}[theorem]{Conjecture}
\newtheorem*{theoremA}{Theorem A}
\newtheorem*{theoremB}{Theorem B}
\newtheorem*{theoremC}{Theorem C}
\newtheorem*{theoremD}{Theorem D}
\newtheorem*{theoremE}{Theorem E}
\newtheorem*{theoremF}{Theorem F}
\newtheorem*{wowconjecture}{Conjecture (WOW-284)}
\theoremstyle{definition}
\newtheorem{definition}[theorem]{Definition}
\theoremstyle{remark}
\newtheorem{remark}[theorem]{Remark}

\BeforeBeginEnvironment{theorem}{\Needspace{6\baselineskip}}
\BeforeBeginEnvironment{proposition}{\Needspace{6\baselineskip}}
\BeforeBeginEnvironment{lemma}{\Needspace{5\baselineskip}}
\BeforeBeginEnvironment{corollary}{\Needspace{5\baselineskip}}
\BeforeBeginEnvironment{proof}{\Needspace{4\baselineskip}}

\newcommand{\F}{\mathbb F_5}
\newcommand{\one}{\mathbf 1}
\newcommand{\Spec}{\operatorname{Spec}}
\newcommand{\tr}{\operatorname{tr}}
\newcommand{\diam}{\operatorname{diam}}
\newcommand{\dist}{\operatorname{dist}}
\newcommand{\lammin}{\lambda_{\min}}
\newcommand{\repoRef}{v2.2.0}
\newcommand{\codefile}[1]{%
  \href{https://github.com/SamPetkov/wow284/blob/\repoRef/#1}{\path{#1}}}

\hypersetup{
  pdftitle={Counterexamples, Spectral Obstructions, and Deletion Stability for WOW-284},
  pdfauthor={Samuil Petkov},
  pdfsubject={Counterexamples and structural results for the minimum-dual-degree and least-distance-eigenvalue conjecture WOW-284},
  pdfkeywords={distance spectrum, dual degree, Moore graph, spectral obstruction, deletion stability},
  pdfcreator={LaTeX with amsart},
  pdfproducer={pdfTeX},
  pdfdisplaydoctitle=true
}

\title[Counterexamples and obstructions for WOW-284]
{Counterexamples, Spectral Obstructions, and Deletion Stability for WOW-284}
\author{Samuil Petkov}
\address{Department of Physics, \'Ecole normale sup\'erieure, Universit\'e PSL, Paris, France}
\email{samuil.petkov@phys.ens.psl.eu}
\date{}
\subjclass[2020]{Primary 05C50; Secondary 05C12, 05E30, 90C22}
\keywords{distance spectrum, dual degree, Moore graph, spectral obstruction, deletion stability}

\begin{document}

\begin{abstract}
WOW-284 predicts that the minimum dual degree of every connected graph of
order at least three and girth at least five is at most the negative of its
least distance eigenvalue.  We first give exact counterexamples of orders
\(38,39,40,42\), and \(50\), including the degree-seven Hoffman--Singleton
graph.  We then study the mechanism of failure.  For regular graphs of
diameter three, the counterexample score is an exact shifted adjacency-spectral
quantity.  Every regular strict counterexample has degree at least six and
diameter at most four; the diameter-four case starts, if at all, in degree ten.
We determine the exact ceiling of the standard one-variable nonbacktracking
linear-programming method and prove that its optimizer is unique up to scale.
Localizing the same polynomial to an edge yields a five-cycle obstruction and
shows that every 6-regular strict counterexample has at most 50 vertices.  We
also derive exact distance spectra for one- and two-vertex punctures of Moore
graphs, prove a general deletion-stability inequality, and determine that the
Hoffman--Singleton counterexample survives every deletion of at most five
vertices, with this radius sharp.  All finite spectral decisions are certified
by exact arithmetic; the principal computer-assisted steps are linked inline
to independent verification scripts.
\end{abstract}

\maketitle

\input{sections/01_introduction}
\input{sections/02_moore_examples}
\input{sections/03_score_calculus}
\input{sections/04_regular_obstructions}
\input{sections/05_lp_edge_local}
\input{sections/06_order50}
\input{sections/07_moore_punctures}
\input{sections/08_deletion_stability}
\input{sections/09_equality_and_obstructions}
\input{sections/10_verification_and_outlook}

\appendix
\input{sections/A_explicit_data}
\input{sections/B_nonadjacent_factorization}

\clearpage
\begingroup
\footnotesize
\setlength{\bibsep}{0pt}
\renewcommand{\bibsection}{\section*{References}}
\bibliographystyle{unsrtnat}
\bibliography{references}
\endgroup

\end{document}
